\documentclass[11pt]{article}

% ── Packages ──────────────────────────────────────────────────────────────────
\usepackage[T1]{fontenc}
\usepackage[utf8]{inputenc}
\usepackage{lmodern}
\usepackage{microtype}
\usepackage{amsmath, amssymb}
\usepackage{booktabs}
\usepackage{array}
\usepackage{xcolor}
\usepackage{colortbl}
\usepackage{geometry}
\usepackage{titlesec}
\usepackage{enumitem}
\usepackage{parskip}
\usepackage{fancyhdr}
\usepackage{hyperref}

% ── Page layout ───────────────────────────────────────────────────────────────
\geometry{
  letterpaper,
  top    = 1in,
  bottom = 1in,
  left   = 1.1in,
  right  = 1.1in
}

% ── Colors ────────────────────────────────────────────────────────────────────
\definecolor{headerblue}{RGB}{31,  73, 125}
\definecolor{rowblue}   {RGB}{220, 230, 241}
\definecolor{ruleblue}  {RGB}{31,  73, 125}
\definecolor{notesgray} {RGB}{100, 100, 100}

% ── Section heading styles ────────────────────────────────────────────────────
\titleformat{\section}
  {\large\bfseries\color{headerblue}}
  {\thesection.}{0.6em}{}[\vspace{-4pt}\textcolor{ruleblue}{\rule{\linewidth}{0.8pt}}]

\titleformat{\subsection}
  {\normalsize\bfseries\color{headerblue}}
  {\thesubsection}{0.5em}{}

\titlespacing*{\section}      {0pt}{18pt}{6pt}
\titlespacing*{\subsection}   {0pt}{12pt}{4pt}

% ── Header / footer ───────────────────────────────────────────────────────────
\pagestyle{fancy}
\fancyhf{}
\renewcommand{\headrulewidth}{0.4pt}
\fancyhead[L]{\small\textcolor{notesgray}{GNN2 Project --- Speed Benchmark}}
\fancyhead[R]{\small\textcolor{notesgray}{March 2026}}
\fancyfoot[C]{\small\thepage}

% ── List spacing ──────────────────────────────────────────────────────────────
\setlist[itemize]{itemsep=3pt, topsep=4pt, leftmargin=1.4em}
\setlist[description]{itemsep=4pt, topsep=4pt, leftmargin=1.6em, labelwidth=0pt, leftmargin=!}

% ── Hyperref ──────────────────────────────────────────────────────────────────
\hypersetup{
  colorlinks = true,
  linkcolor  = headerblue,
  urlcolor   = headerblue,
  pdftitle   = {24-Hour Speed Benchmark: OpenDSS vs.\ GNN Models},
  pdfauthor  = {GNN2 Project},
}

% ─────────────────────────────────────────────────────────────────────────────
\begin{document}

% ── Title block ───────────────────────────────────────────────────────────────
\begin{center}
  {\LARGE\bfseries\color{headerblue}
    24-Hour Speed Benchmark: OpenDSS vs.\ GNN Models}\\[6pt]
  {\large GNN2 Project \quad|\quad March 2026}
\end{center}

\vspace{4pt}
\noindent\textcolor{ruleblue}{\rule{\linewidth}{1.2pt}}
\vspace{10pt}

\tableofcontents
\vspace{14pt}
\noindent\textcolor{ruleblue}{\rule{\linewidth}{0.6pt}}
\vspace{10pt}

% ─────────────────────────────────────────────────────────────────────────────
\section{Executive Summary}
% ─────────────────────────────────────────────────────────────────────────────

This report benchmarks the 24-hour (288 timestep) inference speed of two Graph Neural
Network (GNN) models against OpenDSS, a traditional power-flow solver, on an 89-node
distribution feeder.

\medskip
\noindent\textbf{Core finding:} On CPU, OpenDSS is roughly $22\times$ faster per timestep
than GNN inference.  On GPU with batching, the GNN pipeline comes within $1.6\times$ of
OpenDSS, closing the gap substantially.  The GNN's advantage lies in high-throughput,
batched, parallel scenario evaluation---not in single-scenario latency.

% ─────────────────────────────────────────────────────────────────────────────
\section{Experimental Setup}
% ─────────────────────────────────────────────────────────────────────────────

\subsection{Environment}

\begin{itemize}
  \item \textbf{Working directory:} \verb|C:\Users\alita\OneDrive\Desktop\GNN2|
  \item \textbf{Benchmark script:} \verb|speed_benchmark_and_timing.py|
  \item \textbf{Timesteps:} $N_{\text{steps}} = 288$ \; (24\,h at 5-minute resolution)
  \item \textbf{Network size:} $N = 89$ nodes (training subset)
  \item \textbf{Devices:} CPU and NVIDIA L4 GPU
\end{itemize}

\subsection{Models Evaluated}

\begin{itemize}
  \item \textbf{Loadtype GNN:} \verb|gnn2_architecture_search/loadtype/best.pt|
  \item \textbf{Original GNN:} \verb|gnn2_architecture_search/original/best.pt|
\end{itemize}

\subsection{Inference Modes}

\begin{itemize}
  \item \textbf{Non-batched:} one GNN forward pass per timestep ($288$ separate forwards).
  \item \textbf{Batched:} all $288$ timesteps packed into a single batched forward.
\end{itemize}

OpenDSS runs timestep-by-timestep (no batching is available).  All GNN pipeline times
include feature construction, tensor/batch creation, and the model forward.

% ─────────────────────────────────────────────────────────────────────────────
\section{Notation and Definitions}
% ─────────────────────────────────────────────────────────────────────────────

\subsection{OpenDSS Aggregates}

\[
T^{\text{wall}}_{\text{DSS}} \;\text{(total wall-clock time, 24\,h)},
\qquad
T^{\text{solve}}_{\text{DSS}} \;\text{(sum of pure \texttt{Solve()} times)}.
\]
Per-step means:
\[
\bar{t}^{\text{wall}}_{\text{DSS}}
  = \frac{T^{\text{wall}}_{\text{DSS}}}{N_{\text{steps}}},
\qquad
\bar{t}^{\text{solve}}_{\text{DSS}}
  = \frac{T^{\text{solve}}_{\text{DSS}}}{N_{\text{steps}}}.
\]

\subsection{GNN Aggregates}

\[
T^{\text{fwd}}_{\text{GNN}}
  = \sum_{t=1}^{N_{\text{steps}}} t_{\text{fwd}}(t),
\qquad
\bar{t}^{\text{fwd}}_{\text{GNN}}
  = \frac{T^{\text{fwd}}_{\text{GNN}}}{N_{\text{steps}}}.
\]

\subsection{Solver Iteration Means}

\[
\overline{\text{ControlIter}}
  = \frac{1}{N_{\text{conv}}}
    \sum_{t \in \mathcal{C}} \text{ControlIter}(t),
\qquad
\overline{\text{PFIter}}
  = \frac{1}{N_{\text{conv}}}
    \sum_{t \in \mathcal{C}} \text{PFIter}(t),
\]
where $\mathcal{C}$ is the set of converged steps and $N_{\text{conv}} = |\mathcal{C}|$.

% ─────────────────────────────────────────────────────────────────────────────
\section{Benchmark Results}
% ─────────────────────────────────────────────────────────────────────────────

\subsection{OpenDSS Baseline (CPU)}

\[
T^{\text{wall}}_{\text{DSS}} \approx 0.1763\,\text{s},
\qquad
\bar{t}^{\text{wall}}_{\text{DSS}} \approx 0.612\,\text{ms/step},
\]
\[
T^{\text{solve}}_{\text{DSS}} \approx 0.0253\,\text{s},
\qquad
\bar{t}^{\text{solve}}_{\text{DSS}} \approx 0.088\,\text{ms/step},
\]
\[
\overline{\text{ControlIter}} \approx 2.33,
\qquad
\overline{\text{PFIter}} \approx 4.76.
\]

OpenDSS solves the full 24-hour profile in well under a quarter of a second, with
sub-millisecond solves and only a few nonlinear iterations per step.

\subsection{GNN Forward Times (CPU)}

\paragraph{Loadtype model.}
\[
T^{\text{fwd}}_{\text{GNN,LT}} \approx 3.8749\,\text{s},
\qquad
\bar{t}^{\text{fwd}}_{\text{GNN,LT}} \approx 13.455\,\text{ms/step}.
\]

\paragraph{Original model.}
\[
T^{\text{fwd}}_{\text{GNN,Orig}} \approx 3.8597\,\text{s},
\qquad
\bar{t}^{\text{fwd}}_{\text{GNN,Orig}} \approx 13.402\,\text{ms/step}.
\]

The per-step slowdown relative to OpenDSS wall time is:
\[
\frac{\bar{t}^{\text{fwd}}_{\text{GNN}}}{\bar{t}^{\text{wall}}_{\text{DSS}}}
\approx \frac{13.4\,\text{ms}}{0.612\,\text{ms}} \approx 21.9\times.
\]

\subsection{Summary Table}

\begin{table}[ht]
\centering
\renewcommand{\arraystretch}{1.25}
\begin{tabular}{
  >{\raggedright\arraybackslash}p{5.4cm}
  >{\centering\arraybackslash}p{2.2cm}
  >{\centering\arraybackslash}p{2.4cm}
  >{\centering\arraybackslash}p{2.4cm}
}
\toprule
\rowcolor{headerblue}
\textcolor{white}{\textbf{Configuration}}
  & \textcolor{white}{\textbf{Total (24\,h)}}
  & \textcolor{white}{\textbf{Per-step mean}}
  & \textcolor{white}{\textbf{vs.\ OpenDSS}} \\
\midrule
\rowcolor{rowblue}
OpenDSS (CPU, wall)                  & 0.176\,s   & 0.612\,ms & 1.0$\times$ (baseline) \\
OpenDSS (CPU, solve only)            & 0.025\,s   & 0.088\,ms & --- \\
\rowcolor{rowblue}
GNN Loadtype (CPU, non-batched)      & 3.875\,s   & 13.455\,ms & ${\approx}21.9\times$ slower \\
GNN Original (CPU, non-batched)      & 3.860\,s   & 13.402\,ms & ${\approx}21.9\times$ slower \\
\rowcolor{rowblue}
GNN Original (GPU, non-batched)$^*$  & 0.686\,s   & 2.382\,ms  & ${\approx}1.96\times$ slower \\
GNN Original (GPU, batched)$^*$      & 0.498\,s   & 1.729\,ms  & ${\approx}1.60\times$ slower \\
\bottomrule
\end{tabular}
\caption{Headline benchmark numbers across all configurations.
  $^*$GPU times are from the pipeline breakdown totals (Section~\ref{sec:breakdown}).}
\label{tab:summary}
\end{table}

% ─────────────────────────────────────────────────────────────────────────────
\section{Detailed Pipeline Timing Breakdown}
\label{sec:breakdown}
% ─────────────────────────────────────────────────────────────────────────────

The script separates two pipelines over all 288 steps:
\[
T^{\text{profile}}_{\text{DSS}}
  \approx T_{\text{set\_time}}
         + T_{\text{apply}}
         + T_{\text{solve}}
         + T_{\text{get\_voltage}},
\]
\[
T^{\text{pipe}}_{\text{GNN}}
  = T_{\text{build\_x}}
  + T_{\text{tensor\_creation}}
  + T_{\text{forward}}.
\]

\begin{table}[ht]
\centering
\renewcommand{\arraystretch}{1.25}
\begin{tabular}{
  >{\raggedright\arraybackslash}p{3.6cm}
  >{\centering\arraybackslash}p{1.2cm}
  >{\centering\arraybackslash}p{2.4cm}
  >{\centering\arraybackslash}p{2.4cm}
  >{\centering\arraybackslash}p{1.8cm}
}
\toprule
\rowcolor{headerblue}
\textcolor{white}{\textbf{Dataset / Mode}}
  & \textcolor{white}{\textbf{Device}}
  & \textcolor{white}{\textbf{DSS Profile}}
  & \textcolor{white}{\textbf{GNN Pipeline}}
  & \textcolor{white}{\textbf{Ratio}} \\
\midrule
\rowcolor{rowblue}
Original, non-batched & CPU & 0.396\,s & 4.642\,s & $11.72\times$ \\
Original, non-batched & GPU & 0.349\,s & 0.686\,s & $1.96\times$  \\
\rowcolor{rowblue}
Original, batched     & CPU & 0.315\,s & 3.664\,s & $11.63\times$ \\
Original, batched     & GPU & 0.311\,s & 0.498\,s & $1.60\times$  \\
\bottomrule
\end{tabular}
\caption{Component-level pipeline totals over 288 timesteps.
  Divide by 288 for per-step values.}
\label{tab:breakdown}
\end{table}

\medskip
The main qualitative transitions are:
\begin{itemize}
  \item \textbf{CPU $\to$ GPU:} model forward time drops dramatically, exposing
        feature-construction and tensor-creation overheads as the new bottleneck.
  \item \textbf{Non-batched $\to$ batched:} packing all 288 timesteps into one call
        improves GPU occupancy and amortises kernel-launch overhead, yielding a further
        $\approx 27\%$ speedup on GPU.
\end{itemize}

% ─────────────────────────────────────────────────────────────────────────────
\section{Key Insights and Interpretation}
% ─────────────────────────────────────────────────────────────────────────────

\begin{description}[style=unboxed, leftmargin=0pt, labelwidth=0pt]

  \item[\textbf{Why OpenDSS is fast on this case.}]
  OpenDSS uses a sparse admittance-matrix solver written in optimised C/C\texttt{++}.
  For an 89-node feeder the linear system is tiny and very few nonlinear iterations
  are required per step.  The dominant wall-clock cost is Python API overhead
  (snapshot apply and voltage reads), not the numerical solve itself:
  $\bar{t}^{\text{solve}}_{\text{DSS}} \approx 0.088\,\text{ms}$, while
  $\bar{t}^{\text{wall}}_{\text{DSS}} \approx 0.612\,\text{ms}$.

  \item[\textbf{Why GNN is slower on CPU.}]
  GNN inference involves repeated dense matrix multiplications of size $h \times h$
  across $L$ message-passing layers.  The four primary speed levers are:
  input dimension $d_{\text{node}}$, embedding sizes $(n_{\text{emb}},\,e_{\text{emb}})$,
  hidden width $h$, and depth $L$.  For a small 89-node feeder, this dense computation
  is unnecessarily heavy compared to OpenDSS's sparse approach.

  \item[\textbf{Why GPU and batching close the gap.}]
  GPU Streaming Multiprocessors, warp-level SIMT execution, and Tensor Cores accelerate
  the dense matrix multiplies that dominate GNN inference.  Batching stacks multiple
  timesteps into larger tensor dimensions, improving SM occupancy and amortising
  kernel-launch overhead.  OpenDSS cannot be trivially batched because its iterative
  solver has internal branching and a single-scenario API.

  \item[\textbf{Feature-building overhead.}]
  Once the model forward is accelerated on GPU, the remaining bottleneck shifts to
  Python-side feature construction and tensor/batch creation.  Vectorising feature
  building, eliminating Python loops, and minimising host\,$\leftrightarrow$\,device
  transfers are the next optimisation levers.

  \item[\textbf{Latency vs.\ throughput.}]
  For single-scenario latency, OpenDSS wins on this small feeder.  For high-throughput
  what-if analysis (many scenarios in parallel on GPU), a batched GNN can deliver
  better aggregate throughput.  The present test is a worst-case regime for the GNN:
  small feeder, single scenario, no batching on the baseline.

  \item[\textbf{When does a GNN outperform Newton--Raphson power flow?}]
  For large systems (hundreds to thousands of buses), huge scenario batches, and GPU
  inference, a learned surrogate can match or beat traditional power flow in wall time.
  For $N = 89$ nodes, NR-based power flow is faster due to the small problem dimension
  and the highly efficient sparse solver.  The claim ``GNN is faster than NR'' is
  \emph{not universally true}---it depends on system size, hardware, and batch size.

  \item[\textbf{Two levels of GPU parallelism.}]
  \begin{itemize}
    \item \emph{Within one forward pass (matrix-multiply parallelism):}
          GPU SIMT execution and Tensor Cores (FP16/BF16/TF32) accelerate dense
          operations, reducing $T_{\text{forward}}$ dramatically.
    \item \emph{Across scenarios/timesteps (batch parallelism):}
          batching creates larger tensor dimensions, giving the GPU more work per
          kernel launch and improving occupancy across SMs.
  \end{itemize}

  \item[\textbf{MLP vs.\ GNN trade-off.}]
  A smaller, shallower MLP surrogate could potentially close the remaining gap even for
  89 nodes, trading some accuracy for speed.  The current GNN architecture
  (message passing + dense MLP blocks) may be heavier than necessary for a small feeder.

  \item[\textbf{Consistency with the literature.}]
  Claims that learned surrogates are faster than traditional power flow typically assume
  larger networks (hundreds--thousands of buses), GPU inference, and/or many scenarios
  evaluated in parallel.  In the present setting
  ($N = 89$, Python overhead, single-scenario evaluation), it is expected that OpenDSS
  wins on CPU, while GPU and batching reduce the gap substantially.

\end{description}

% ─────────────────────────────────────────────────────────────────────────────
\section{Conclusion and Next Steps}
% ─────────────────────────────────────────────────────────────────────────────

On the 89-node feeder tested, OpenDSS (CPU) is the fastest solver for single-scenario
latency, completing 288 timesteps in under 0.18\,s.  GNN inference on CPU is
${\approx}22\times$ slower per step, but GPU acceleration and batching reduce this gap
to just ${\approx}1.6\times$.  These results are consistent with the broader literature:
GNN surrogates offer their greatest advantage at scale rather than for small,
single-scenario evaluations.

\medskip
\noindent\textbf{Priority next steps:}
\begin{enumerate}[itemsep=3pt]
  \item \textbf{Reduce Python feature-building overhead} --- vectorise construction loops
        and minimise host\,$\leftrightarrow$\,device transfers.
  \item \textbf{Explore lighter MLP architectures} for small feeders where full
        message-passing GNNs may be over-parameterised.
  \item \textbf{Scale tests to larger networks} (hundreds--thousands of buses) where GNN
        parallelism becomes the decisive advantage.
  \item \textbf{Profile batch-size sensitivity} to identify the optimal batch size for
        the GPU configuration in use.
\end{enumerate}

\end{document}
